% =====================================================================
% Chapter 1 -- Introduction
% =====================================================================
\documentclass[../preamble.tex]{subfiles}
\begin{document}

\chapter{Introduction}
\label{ch:introduction}

\section{Clinical context}
\label{sec:intro-context}

Polycystic Ovary Syndrome (PCOS) is a heterogeneous endocrine disorder that
affects reproductive, metabolic, and psychological health. Its assessment is
not based on a single diagnostic test. The Rotterdam consensus describes a
combination of oligo- or anovulation, clinical or biochemical
hyperandrogenism, and polycystic ovarian morphology, with two of the three
features required after alternative causes have been excluded
\cite{rotterdam2003}. International guidelines also emphasise that diagnosis
requires clinical context rather than image classification alone
\cite{teede2018}.

Ultrasound is important in this process, but ovarian ultrasound images are
not uniform measurements. Speckle noise, low contrast, probe position,
operator technique, acquisition device, compression, and the distribution of
visible follicles all change the appearance of an image. A model can therefore
obtain a high score on a single source while learning properties that do not
survive a dataset change. This is especially concerning when a false positive
or false negative changes the clinical pathway.

\section{Research problem}
\label{sec:intro-problem}

Many PCOS image-classification studies report accuracy or AUC on a single
public dataset. Such numbers are useful for comparing models under a fixed
protocol, but they do not answer three practical questions:

\begin{enumerate}
  \item Does a transfer-learned model generalise to a separate PCOS ultrasound
        cohort?
  \item Does a probability such as $0.95$ correspond to a reliable confidence
        estimate, or is the model overconfident?
  \item Can a clinician inspect why the model made a prediction and identify
        cases that should be referred for human review?
\end{enumerate}

This thesis addresses these questions through \PEARL{}, a pipeline that
combines transfer learning, ultrasound-specific preprocessing, cross-dataset
evaluation, temperature calibration, MC-Dropout uncertainty estimation, and
Grad-CAM visual explanation. The project does not claim that an image
classifier replaces clinical diagnosis. Its aim is to evaluate whether a
model can be made more transparent about its own performance and its own
limitations.

\section{Research questions}
\label{sec:intro-rq}

The original project proposal defined three broad research questions. They
were refined as the experiments moved from a single-dataset study to an
external-validation study:

\begin{description}
  \item[RQ1.] How do nine transfer-learned backbones respond to different
    ultrasound preprocessing choices, and does the ranking observed on the
    source dataset survive evaluation on PCOSgen?
  \item[RQ2.] Does fine-tuning on a target-cohort training pool recover the
    cross-dataset performance lost by foundation checkpoints, and do
    calibration and MC-Dropout provide useful information beyond a single
    softmax score?
  \item[RQ3.] Do Grad-CAM explanations expose visually plausible attention and
    reveal failure modes or shortcut behaviour that aggregate metrics hide?
\end{description}

The proposal initially described a six-condition preprocessing grid, multiple
XAI methods, and an uncertainty-gated referral system. The executed scope was
later narrowed to a controlled two-denoiser comparison, Grad-CAM, and a
referral analysis as a design concept rather than a deployed clinical system.
The evolution is recorded explicitly in \cref{app:scope-evolution} rather than
silently presenting the original proposal as completed work.

\section{Objectives}
\label{sec:intro-objectives}

The executed objectives were:

\begin{enumerate}
  \item deduplicate and split the Figshare source data, then benchmark nine
        standard transfer-learning backbones;
  \item compare SRAD and Gaussian denoising while holding the other
        preprocessing steps fixed;
  \item evaluate zero-shot generalisation on the PCOSgen corpus;
  \item fine-tune selected models on a PCOSgen training pool while preserving
        a held-out test set;
  \item construct a probability-averaged three-model ensemble;
  \item measure discrimination, threshold behaviour, calibration, and
        predictive entropy; and
  \item use Grad-CAM to inspect representative correct, uncertain, and
        failure cases.
\end{enumerate}

\section{Contributions}
\label{sec:intro-contributions}

This thesis makes the following contributions:

\begin{enumerate}
  \item It documents a reproducible PCOS ultrasound pipeline that separates
        source-domain training, target-domain adaptation, calibration,
        uncertainty estimation, and explanation.
  \item It reports a controlled comparison of SRAD and Gaussian denoising and
        shows that letterbox padding can be harmful under cross-dataset
        evaluation.
  \item It demonstrates the difference between strong internal performance
        and poor zero-shot external performance, then measures recovery after
        target-cohort fine-tuning.
  \item It reports a pre-HPO probability-averaged ensemble with AUC 0.9487,
        F1 0.8665, MCC 0.7884, Brier score 0.0981, and calibrated ECE 0.0810
        on the held-out PCOSgen test set.
  \item It combines MC-Dropout entropy and Grad-CAM inspection with the
        performance analysis, while clearly distinguishing descriptive
        uncertainty analysis from a clinically validated referral rule.
  \item It records the differences between the intended v3 pipeline design
        and the executed final run so that future work can reproduce or extend
        the missing EMA, SWA, k-fold, LRP, and SHAP experiments.
\end{enumerate}

\section{Thesis organisation}
\label{sec:intro-organisation}

\Cref{ch:background} reviews the clinical and technical background.
\Cref{ch:method-overview} presents the complete pipeline. The datasets and
splits are described in \cref{ch:data}, followed by preprocessing in
\cref{ch:preprocessing} and model training in \cref{ch:training}. Calibration,
uncertainty, and XAI are described in \cref{ch:trust}. \Cref{ch:results}
reports the experiments, \cref{ch:discussion} interprets their implications,
and \cref{ch:conclusion} concludes the work. The appendices preserve the
intended-but-unexecuted design, HPO details, scope changes, reproducibility
instructions, and the full intended-versus-executed comparison.

\end{document}
